\section{Setup and assumptions}

Generic positive constants may change between displays; the model constants introduced below remain fixed across laws and sample sizes. For a vector, \(\|\cdot\|_\infty\) denotes the maximum absolute coordinate, and for a function it denotes the supremum of its absolute value over its stated domain. Multi-indices have nonnegative integer coordinates, with \(|\alpha|\) denoting their sum. Observation indices range from \(1\) to \(n\), and \(a\vee b=\max\{a,b\}\). The notation \(\lesssim\) denotes an inequality up to a positive constant depending on the fixed model parameters, and \(\asymp\) denotes inequalities in both directions. For comparisons uniform over the overlap range specified below, these constants are independent of both the sample size and the overlap exponent.

\subsection{Sample and target}

Consider an \leanref{S-2}{observational causal model} indexed by a probability law \leanref{sym:P}{\ensuremath{P}}. The covariate vector \leanref{sym:X}{\ensuremath{X}} takes values in the covariate cube \leanref{sym:mathcalX}{\ensuremath{\mathcal X}}, where \(\mathcal X=[0,1]^d\) and \(d\ge1\). The binary treatment indicator \leanref{sym:A}{\ensuremath{A}} takes values in \(\{0,1\}\). Write \(Y(0)\) for the control potential outcome, \leanref{sym:Y1}{\ensuremath{Y(1)}} for the treated potential outcome, and \leanref{sym:Y}{\ensuremath{Y}} for the observed outcome. The observed triple \leanref{sym:Z}{\ensuremath{Z}} is \(Z=(X,A,Y)\). When specifying the sampling law, \(P\) denotes the observed-data marginal of the causal law, and \(P_X\) denotes its covariate marginal.

The covariate density \leanref{sym:f}{\ensuremath{f}} is the density of \(X\) with respect to Lebesgue measure on \(\mathcal X\), and the propensity \leanref{sym:e}{\ensuremath{e(x)}} is a measurable version of \(P(A=1\mid X=x)\). Our target is the treated causal response \leanref{sym:mu1}{\ensuremath{\mu_1}}, the continuous Hölder representative of
\[
\mu_1(x)=\mathbb E_P[Y(1)\mid X=x].
\]
We write \(\mu_{1,P}\) when the indexing law needs to be explicit.

The \leanref{S-3}{sample and estimator convention} uses the observed sample \leanref{sym:Dn}{\ensuremath{\mathcal D_n}}, defined by \(\mathcal D_n=(Z_i)_{i=1}^n\), where \(Z_i=(X_i,A_i,Y_i)\). Every estimator is measurable with respect to this sample. Scalar-valued estimators estimate the response at a specified covariate point; function-valued estimators recover the response over \(\mathcal X\) and are required to have measurable spatial-supremum loss. Expectations of their losses refer to the observed-sample distribution.

Fix a smoothness order \(\beta>1\), a Hölder radius \leanref{sym:L}{\ensuremath{L}} with \(L>0\), and a treated-response bound \leanref{sym:B}{\ensuremath{B}} with \(B>0\). Also fix a density lower bound \leanref{sym:c_f}{\ensuremath{c_f}} with \(0<c_f\le1\) and a global-tail constant \leanref{sym:C}{\ensuremath{C}} with \(C\ge1\). The overlap exponent ranges over the compact interval \leanref{sym:Gamma}{\ensuremath{\Gamma}}, defined by
\[
\Gamma=[\gamma_{\min},\gamma_{\max}].
\]
Its lower endpoint \leanref{sym:gamma_min}{\ensuremath{\gamma_{\min}}} and upper endpoint \leanref{sym:gamma_max}{\ensuremath{\gamma_{\max}}} satisfy \(1<\gamma_{\min}\le\gamma_{\max}<\infty\). Throughout, \(\gamma\in\Gamma\), and \(n\ge1\).

\subsection{Intrinsic cube smoothness}

Smoothness is measured directly on the closed covariate cube. For the fixed smoothness order, the derivative order \leanref{sym:m}{\ensuremath{m}} is \(m=\lceil\beta\rceil-1\). For a generic function \leanref{sym:g}{\ensuremath{g}} on \(\mathcal X\), \(D^\alpha g\) denotes its coordinate derivative, interpreted through continuous boundary traces as specified below. The intrinsic cube Hölder ball \leanref{sym:HbetaL}{\ensuremath{\mathcal H^\beta(L)}} combines a uniform derivative bound with a modulus restriction on the highest-order derivatives. We first state these two restrictions.

\begin{assumptionv}{ass:holder-derivative-bound}[Uniform bound on derivatives]
\[
\max_{|\alpha|\le m}\|D^\alpha g\|_\infty\le L.
\]
\end{assumptionv}

The uniform derivative bound controls the function itself and all coordinate derivatives through order \(m\), using a common radius. This standard Hölder-ball condition gives a common scale for response levels and local polynomial coefficients \citep{Dorn2026GlobalOverlap}. The corresponding restriction on variation across covariate points is as follows.

\begin{assumptionv}{ass:holder-modulus}[Hölder modulus of derivatives]
\[
\max_{|\alpha|=m}\sup_{x\ne x'}
\frac{|D^\alpha g(x)-D^\alpha g(x')|}
{\|x-x'\|_\infty^{\beta-m}}
\le L.
\]
\end{assumptionv}

The Hölder modulus condition bounds changes in the highest-order derivatives by the covariate distance raised to \(\beta-m\). Together with the derivative bound, this standard restriction controls local approximation error at smoothness order \(\beta\) \citep{Dorn2026GlobalOverlap}. The next definition fixes the boundary convention used throughout.

\begin{definitionv}{def:holder}[Intrinsic cube Hölder ball]
The intrinsic cube Hölder ball is
\[
\mathcal H^\beta(L)
:=
\left\{
g:\mathcal X\to\mathbb R:
g\in C^m(\mathcal X),\
g\text{ satisfies }\cref{obj:ass:holder-derivative-bound,obj:ass:holder-modulus}
\right\},
\qquad
\mathcal X=[0,1]^d,
\qquad
m=\max\{\lceil\beta\rceil-1,0\}.
\]
Here \(d\) is a nonnegative integer and \(\beta,L\in\mathbb R\). The integer \(m\) is the derivative order, and the top-order modulus has exponent \(\beta-m\). The class \(C^m(\mathcal X)\) consists of functions whose classical coordinate partial derivatives through order \(m\) on the interior of the cube admit unique continuous traces on \(\mathcal X\). Every derivative \(D^\alpha g\), supremum, and difference quotient in the defining restrictions uses these traces. At positive integer \(\beta\), this convention uses Lipschitz derivatives of order \(\beta-1\).
\end{definitionv}

For the maintained domain \(\beta>1\), the derivative order in \cref{obj:def:holder} agrees with the order fixed above. Its trace convention makes smoothness restrictions meaningful at boundary points as well as in the interior.

\subsection{Sampling and observational restrictions}

The statistical experiment combines independent sampling with restrictions on covariate coverage, causal identification, treatment availability, and the treated response. We begin with the sampling law.

\begin{assumptionv}{ass:iid}[Independent and identically distributed observations]
The observed sample \(\mathcal D_n\) consists of \(n\) independent observations with common law \(P\), so that
\[
\mathcal D_n\sim P^{\otimes n}.
\]
\end{assumptionv}

Independent and identically distributed sampling gives the product experiment \(P^{\otimes n}\) used to evaluate estimation risk. It is a standard sampling condition for nonparametric regression \citep{Stone1982}. Covariate coverage is controlled separately.

\begin{assumptionv}{ass:density}[Covariate density lower bound]
The covariate density \(f\) satisfies
\[
c_f\le f(x)
\]
for Lebesgue-almost every \(x\in\mathcal X\).
\end{assumptionv}

The random-design density lower bound ensures that each measurable covariate region has probability at least \(c_f\) times its volume. This standard coverage condition supports response recovery throughout the cube \citep{Stone1982}. To connect that response to observed treated outcomes, we impose consistency and conditional exchangeability.

\begin{assumptionv}{ass:consistency}[Consistency of observed outcomes]
The observed outcome \(Y\), binary treatment indicator \(A\), and potential outcomes \(Y(0)\) and \(Y(1)\) satisfy
\[
Y=AY(1)+(1-A)Y(0)
\]
almost surely.
\end{assumptionv}

Consistency links each observed outcome to the potential outcome corresponding to the realized treatment. It is the standard condition that gives the observed treatment groups their potential-outcome interpretation \citep{KennedyBalakrishnanRobinsWasserman2024}. The next condition makes treatment assignment comparable within covariate strata.

\begin{assumptionv}{ass:exchangeability}[Conditional exchangeability of treatment]
The potential outcomes \(Y(0)\) and \(Y(1)\) are jointly independent of the treatment indicator \(A\) conditional on the covariate vector \(X\):
\[
(Y(0),Y(1))\perp A\mid X.
\]
\end{assumptionv}

Conditional exchangeability is the standard identification condition under which conditioning on \(X\) removes dependence between treatment assignment and potential outcomes \citep{KennedyBalakrishnanRobinsWasserman2024}. Together with consistency, it identifies the treated causal response from the treated conditional mean wherever the propensity is positive. Treatment availability is quantified by a global tail restriction.

\begin{assumptionv}{ass:global-tail}[Global propensity tail bound]
The propensity \(e(X)\), evaluated at the covariate vector \(X\), satisfies
\[
P\{e(X)\le t\}\le Ct^{\gamma-1}
\qquad\text{for every }t\in[0,1].
\]
\end{assumptionv}

The polynomial global propensity-tail bound controls the covariate probability assigned to regions with low treatment probability, a standard overlap restriction in this setting \citep{Dorn2026GlobalOverlap}. Since \(\gamma>1\), its value at \(t=0\) implies positive propensity almost surely under \(P_X\). Consequently, consistency and exchangeability identify \(\mu_1(x)\) with \(\mathbb E_P[Y\mid A=1,X=x]\) for \(P_X\)-almost every \(x\). The smoothness convention selects a continuous representative over the entire cube.

We next control the magnitude and fluctuations of the treated conditional outcome distribution.

\begin{assumptionv}{ass:bounded-potential-outcomes}[Treated response magnitude bound]
The observed outcome \(Y\), binary treatment indicator \(A\), and covariate vector \(X\) satisfy
\[
\operatorname*{ess\,sup}_{x\sim P_X}
\left(
 \left|\mathbb E_P[Y\mid A=1,X=x]\right|
 \vee
 \sup_{\lambda\in\mathbb R\setminus\{0\}}
 \frac{\sqrt{2\log \mathbb E_P\!\left[
   \exp\!\left\{\lambda\bigl(Y-\mathbb E_P[Y\mid A=1,X]\bigr)\right\}
   \mid A=1,X=x\right]}}{|\lambda|}
\right)
\le B.
\]
\end{assumptionv}

The uniformly bounded treated conditional mean and conditional sub-Gaussian residual condition supplies a common scale for response levels and sampling fluctuations \citep{Dorn2026GlobalOverlap}. In particular, the treated conditional mean has absolute value at most \(B\), and the centered treated residual has conditional moment-generating function bounded by \(\exp(\lambda^2B^2/2)\). These bounds hold essentially uniformly over the covariate distribution. The response curve also satisfies the following smoothness restriction.

\begin{assumptionv}{ass:holder-response}[Hölder smoothness of response]
The treated causal response \(\mu_1\) belongs to the intrinsic cube Hölder ball \(\mathcal H^\beta(L)\):
\[
\mu_1\in\mathcal H^\beta(L).
\]
\end{assumptionv}

Isotropic Hölder response smoothness imposes the same order \(\beta\) across covariate directions, with the common derivative and modulus bounds specified in \cref{obj:def:holder} \citep{Dorn2026GlobalOverlap}. It supplies the regularity through which nearby treated observations inform the target response. Combined with identification and the density lower bound, continuity extends the treated-mean magnitude bound to \(|\mu_1(x)|\le B\) throughout \(\mathcal X\).

\subsection{Law class and risk scales}

We collect the observational restrictions into a fixed-constant class. The dimension, smoothness order, and constants \(L,B,c_f,C\) remain common across its members; \(\gamma\) indexes the overlap restriction. The observational causal law class \leanref{sym:Pgamma}{\ensuremath{\mathcal P_\gamma}} is defined as follows.

\begin{definitionv}{def:model}[Observational causal law class $\mathcal P_\gamma$]
The class \(\mathcal P_\gamma\) consists of observational causal laws \(P\) satisfying the following restrictions, with common fixed constants:
\begin{itemize}
\item \textbf{(Covariate density.)} The density \(f\) of the covariate vector \(X\) satisfies \cref{obj:ass:density}.
\item \textbf{(Consistency.)} The observed outcome \(Y\), binary treatment indicator \(A\), and potential outcomes \(Y(0)\) and \(Y(1)\) satisfy \cref{obj:ass:consistency}.
\item \textbf{(Exchangeability.)} The potential outcomes and treatment satisfy \cref{obj:ass:exchangeability}.
\item \textbf{(Overlap.)} The propensity \(e(X)\) satisfies \cref{obj:ass:global-tail}.
\item \textbf{(Response bound.)} The treated conditional outcome distribution satisfies \cref{obj:ass:bounded-potential-outcomes}.
\item \textbf{(Response smoothness.)} The treated causal response \(\mu_1\) satisfies \cref{obj:ass:holder-response}.
\end{itemize}
\end{definitionv}

For each member of \cref{obj:def:model}, the sample is drawn according to \cref{obj:ass:iid}. Fix an evaluation point \leanref{sym:x0}{\ensuremath{x_0}} in the closed cube, including its boundary. We measure recovery through the pointwise minimax risk \leanref{sym:Rpt}{\ensuremath{R_{\mathrm{pt}}(n,\gamma,x_0)}} and the spatial-supremum minimax risk \leanref{sym:Rsup}{\ensuremath{R_\infty(n,\gamma)}}. Both criteria use expected absolute error, with the latter evaluating the largest error across the cube within each sample realization.

\begin{definitionv}{def:risks}[Minimax risks $R_{\mathrm{pt}}$ and $R_\infty$]
The minimax expected absolute error at \(x_0\) and the minimax expected spatial-supremum error are, respectively,
\[
R_{\mathrm{pt}}(n,\gamma,x_0)
=
\inf_T\sup_{P\in\mathcal P_\gamma}
\mathbb E_P\left|T-\mu_{1,P}(x_0)\right|
\]
and
\[
R_\infty(n,\gamma)
=
\inf_T\sup_{P\in\mathcal P_\gamma}
\mathbb E_P\left[
\sup_{x\in\mathcal X}\left|T(x)-\mu_{1,P}(x)\right|
\right].
\]
Here \(\mu_{1,P}\) is the treated response under law \(P\), and each infimum ranges over estimators measurable with respect to the observed sample \(\mathcal D_n\). The first infimum uses scalar-valued estimators, and the second uses function-valued estimators.
\end{definitionv}

The order of supremum and expectation in \cref{obj:def:risks} makes the functional criterion sensitive to the largest spatial error in each realized estimate. Its spatial supremum ranges over all points of \(\mathcal X\), using the continuous response representative fixed above.

To express the subsequent risk comparisons, define the overlap-dependent effective dimension \leanref{sym:Dgamma}{\ensuremath{D_\gamma}}, the oracle mesh scale \leanref{sym:hgamman}{\ensuremath{h_{\gamma,n}}}, and the response-recovery scale \leanref{sym:rgamman}{\ensuremath{r_{\gamma,n}}} by
\[
D_\gamma=\frac{d\gamma}{\gamma-1},
\qquad
h_{\gamma,n}=n^{-1/(2\beta+D_\gamma)},
\qquad
r_{\gamma,n}=h_{\gamma,n}^{\beta}
=n^{-\beta/(2\beta+D_\gamma)}.
\]
These deterministic scales depend on the sample size, smoothness order, dimension, and overlap exponent. They satisfy
\[
h_{\gamma,n}^{\beta}
=
\bigl(nh_{\gamma,n}^{D_\gamma}\bigr)^{-1/2},
\]
which expresses the balance between a smoothness scale and an overlap-adjusted sampling scale. The term ``oracle'' indicates that the mesh scale is indexed by \(\gamma\); it provides the benchmark for comparing risk and selecting a mesh from the observed treatment design.
